% =====================================================================
%  Section 06 — Phase 1 spatial / attribute merges A–G
% =====================================================================

\section{Phase 1 \textbullet\ The Seven Merges (A--G)}
\label{sec:phase1-merges}

The cleaning notebook builds seven cross-dataset merges. Each one solves
one analytical problem and is consumed directly by a specific subset of
research questions. They are exported as CSVs to \fn{CleanedData/} so
that the analysis notebook depends only on stable inputs.

\subsection{Merge index}

\rowcolors{2}{darkpanel}{darkbg}
\begin{table}[H]
\centering\small
\begin{tabular}{@{} >{\color{plotorange}\bfseries}l
                    >{\color{bodytext}}l
                    >{\color{dimtext}}l
                    >{\color{bodytext}}l @{}}
\toprule
\color{plotblue}\# & \color{plotblue}Datasets joined & \color{plotblue}Predicate & \color{plotblue}Serves \\
\midrule
A & Boarding $\to$ Population (point-in-polygon) & \fn{within}    & Q1, Q9 \\
B & Terminal centroids $\to$ Population          & \fn{within}    & Q2, Q5 \\
C & Routes $\to$ Terminals                       & attribute join & Q10, Q11 \\
D & Vehicle/Pax Flow $\cap$ Terminal 50-m buffer & \fn{intersects} & Q2, Q10 (G2) \\
E & Merge B + Merge D                            & outer on \fn{Terminal\_ID} & Q2 \\
F & Routes-per-terminal aggregate                & group-by       & Q11 \\
G & Underused-terminal flag (F + 100-m buffer)   & \fn{within}    & Q11 (G1) \\
\bottomrule
\end{tabular}
\caption{Merge index. Predicates are the GeoPandas spatial-relationship
operators used in \fn{gpd.sjoin}.}
\end{table}
\rowcolors{0}{}{}

\subsection{Merge A — Boarding into Population hexagons (Q1, Q9)}

A point-in-polygon spatial join assigns each boarding stop to the H3
hexagon that contains it. Every stop inherits the hex's
\fn{Population\_2018}, \fn{Jobs\_In\_Hex}, and
\fn{Jobs\_Accessible\_60min} columns. This is the foundation of every
question that needs ``how many people / how many jobs'' near a stop.

\begin{lstlisting}[caption={Merge A — point-in-polygon spatial join.}]
merge_a = gpd.sjoin(
    boarding,
    population[[
        "Population_2018", "Pop_Male", "Pop_Female",
        "Jobs_In_Hex", "Jobs_Accessible_60min",
        "Jobs_Pct_Accessible_60min", "H3_Hex_ID", "geometry",
    ]],
    how="left", predicate="within",
)
\end{lstlisting}

\subsection{Merge B — Terminal centroids into Population (Q2, Q5)}

Terminal polygons are reduced to centroids before the spatial join,
so that each terminal gets exactly one hex's worth of demographic
context.

\subsection{Merge C — Routes enriched with terminal names (Q10, Q11)}

Routes carry only numeric terminal IDs. Two left-joins on a terminal
lookup table (one for origin, one for destination) attach human-readable
names so that downstream tables and visualisations can label terminals
directly.

\subsection{Merge D — Flow at terminal buffers (Q2, Q10 / G2)}

The pivotal merge for the \emph{Empty Return Index}. A 50-m buffer is
drawn around each terminal centroid; every passenger-flow segment and
every vehicle-flow segment that intersects the buffer is summed:

\begin{lstlisting}[caption={Merge D — sum vehicle and passenger flow within 50 m of each terminal.}]
BUFFER_RADIUS = 50  # metres
results = []
for _, term in terminals.iterrows():
    buf = term.geometry.centroid.buffer(BUFFER_RADIUS)
    veh_mask = veh_flow.geometry.intersects(buf)
    pax_mask = pax_flow.geometry.intersects(buf)
    total_veh = veh_flow.loc[veh_mask, "Vehicles_Per_Hour"].sum()
    total_pax = pax_flow.loc[pax_mask, "Passengers_Per_Hour"].sum()
    ratio = total_pax / total_veh if total_veh > 0 else np.nan
    results.append({
        "Terminal_ID":            term["Terminal_ID"],
        "Vehicles_At_Terminal":   total_veh,
        "Passengers_At_Terminal": total_pax,
        "Pax_Per_Vehicle_Ratio":  ratio,
    })
merge_d = pd.DataFrame(results)
merge_d["Empty_Return_Index"] = np.where(
    merge_d["Vehicles_At_Terminal"] > 0,
    (1 - merge_d["Passengers_At_Terminal"]
         / merge_d["Vehicles_At_Terminal"]).clip(lower=0),
    np.nan,
)
\end{lstlisting}

\subsection{Merge E — Terminal demographics + flow (Q2)}

A simple outer-join of B and D on \fn{Terminal\_ID}. Each terminal now
has both its surrounding population context and its measured flow
volumes. Used by Q2 to ask whether terminals with high passengers /
vehicle ratio are sited in dense areas.

\subsection{Merge F — Routes per terminal (Q11)}

A \fn{group-by} aggregate on origin and destination separately, then
joined back to the terminal lookup. Each terminal gets:
\fn{Routes\_As\_Origin}, \fn{Routes\_As\_Dest}, \fn{Total\_Routes}, plus
average route lengths in each direction.

\subsection{Merge G — Underused-terminal flag (Q11 / G1)}

Combines F with a 100-m buffer around each terminal centroid; a terminal
is flagged \fn{Underused} when it has at least one route but zero
boarding within the buffer. The notebook also splits terminals into
four operational zones using median-route and median-boarding splits:

\begin{itemize}
  \item \kw{Efficient Hub} — high routes, high boarding.
  \item \kw{Overcrowded / Underserved} — low routes, high boarding.
  \item \kw{Dead Weight} — high routes, low boarding (the policy problem).
  \item \kw{Marginal} — low routes, low boarding.
\end{itemize}

\subsection{Question-specific helper datasets}

Eight smaller helper datasets are also written so that the analysis
notebook does not have to re-derive features per question. Each is
named for its consuming question (\fn{q3\_drop\_points.csv},
\fn{q4\_formal\_vs\_informal.csv}, \fn{q5\_hex\_terminal\_counts.csv},
\fn{q6\_routes\_vs\_passengers.csv}, \fn{q7\_morning\_vs\_evening.csv},
\fn{q8\_spatial\_autocorrelation.csv}, \fn{q9\_underserved.csv},
\fn{q12\_vehicle\_type\_route\_length.csv}). Every one has a
deterministic recipe in the notebook.
